\section{Generic and torsion-specialized monodromy}
\label{sec:generic-monodromy}

The period-$4$ sextic has full symmetric monodromy generically, and
this persists at all but finitely many torsion fibers.  Define
\[
 P_X(T)=T^6+2T^5+4T^4+6T^3+(X-5)T^2-8T-16.
 \numberthis\label{eq:generic-sextic}
\]

\begin{theorem}[Generic monodromy]
\label{thm:generic-S6}
The geometric and arithmetic Galois groups of $P_X$ over
$\overline\Q(X)$ and $\Q(X)$ are both $S_6$.  Its splitting field over
$\Q(X)$ is regular over $\Q$.
\end{theorem}

\begin{proof}
The discriminant is
\[
 \operatorname{disc}_T(P_X)=2^{10}Q(X),
 \numberthis\label{eq:generic-discriminant}
\]
where
\begin{align*}
 Q(X)={}&16X^6-619X^5+23722X^4-366579X^3\\
       &+2799652X^2-12598144X+60081152.
\end{align*}
The polynomial $Q$ is squarefree.  For a short exact certificate,
reduce modulo $5$:
\[
 \overline Q=X^6+X^5+2X^4+X^3+2X^2+X+2,
\]
\[
 \overline Q'=X^5+3X^3+3X^2+4X+1.
\]
The monic nonzero remainders in the Euclidean algorithm are
\[
 X^4+4X+4,\quad X^3+3X^2+2,\quad X^2+3X,\quad1.
\]
Thus $Q$ has six simple roots; its constant term also shows that none
is zero.

The equation $P_X(T)=0$ is the degree-six rational cover
$X=g(T)$ from the period-$4$ parametrization.  Above $X=\infty$ its
two poles have orders $4$ and $2$, so the inertia type is $(4)(2)$ and
its ramification contribution is $4$.  Each of the six finite simple
zeros of the discriminant gives a transposition, contributing $6$.
These contributions total $10=2\cdot6-2$, as required by
Riemann--Hurwitz for a degree-six map from $\mathbf P^1$.

The cover is connected, and its finite branch cycles are
transpositions.  Their product is the inverse of the branch cycle at
infinity, so they generate the full monodromy group.  A transitive
permutation group generated by transpositions is the full symmetric
group: the transpositions are the edges of a connected graph on the
six sheets.  Hence the geometric group is $S_6$.  The arithmetic group
contains the geometric group and is itself a subgroup of $S_6$, so it
is also $S_6$.  Equality of the two groups is equivalent to regularity
of the splitting field.
\end{proof}

\begin{theorem}[Cofinite torsion specialization]
\label{thm:cofinite-torsion-S6}
For all but finitely many roots of unity $\zeta$,
\[
 \Gal(P_\zeta/\Q^{\mathrm{ab}})\cong S_6.
 \numberthis\label{eq:torsion-S6}
\]
For every such nontrivial $\zeta$, if $K=\Q(\zeta)$ and $L$ is the
splitting field of $P_\zeta$ over $K$, then
\[
 \Gal(L/K)=S_6,\qquad
 L\cap K^{\mathrm{ab}}=K\bigl(\sqrt{Q(\zeta)}\bigr),
 \numberthis\label{eq:sign-intersection}
\]
and
\[
 \Gal(LK^{\mathrm{ab}}/K^{\mathrm{ab}})=A_6.
 \numberthis\label{eq:A6-over-Kab}
\]
\end{theorem}

\begin{proof}
Let $\mathcal L/\Q(X)$ be the splitting field of $P_X$.  By
Theorem~\ref{thm:generic-S6}, it is regular and has group $G=S_6$.
After removing the branch locus and the points $0,\infty$, its
normalization gives a finite \'etale $G$-cover over an open subset of
$\mathbf G_m$.  For each conjugacy class of maximal subgroups $M<G$,
the fixed field $\mathcal L^M$ is the function field of a geometrically
integral quotient cover defined over $\Q$; regularity gives
\[
 (\mathcal L^M)\overline\Q=(\mathcal L\overline\Q)^M.
\]
We claim that every one of these finitely many quotient covers satisfies
(PB).

Suppose first that $M\ne A_6$.  If (PB) failed, then
\cite[Proposition 2.1]{Zannier2010}, applied over $\overline\Q$, would
give a factorization through a nontrivial isogeny of $\mathbf G_m$.
Thus, for some $n>1$, the cyclic extension
\[
 \overline\Q(X)(U),\qquad U^n=X,
\]
would be contained in $(\mathcal L\overline\Q)^M$.  If
\[
 N=\Gal(\mathcal L\overline\Q/\overline\Q(U)),
\]
then $M\subseteq N\lhd G$ and $N\ne G$.  Maximality gives $N=M$, so
$M$ is normal and $G/M$ is a nontrivial cyclic quotient.  The only such
maximal kernel in $S_6$ is $A_6$, a contradiction.  Hence this quotient
has (PB).

For $M=A_6$, the quotient is the sign cover
\[
 W^2=Q(X).
\]
For every $n\ge1$, the polynomial $Q(U^n)$ has a simple nonzero zero
and hence is not a square in $\overline\Q(U)$.  Thus
$W^2-Q(U^n)$ is irreducible, so this quotient also has (PB).

Apply Zannier's theorem \cite[Theorem 2.1]{Zannier2010} to all these
quotient covers simultaneously, and enlarge the finite exceptional set
by their branch and bad-specialization points.  For a remaining root of
unity $\zeta$, let $H_\zeta\le G$ be the image of
$\Gal(\overline\Q/\Q^{\mathrm{ab}})$ on a geometric
fiber of the $G$-cover.  It is the Galois group of $P_\zeta$ over
$\Q^{\mathrm{ab}}$.  The fiber of the $M$-quotient is irreducible
precisely when $H_\zeta$ acts transitively on $G/M$.  If $H_\zeta$ were
proper, it would lie in a conjugate of a maximal subgroup $M$ and would
therefore fix a coset in $G/M$.  Since $[G:M]>1$, this contradicts
transitivity.  Hence $H_\zeta=G=S_6$, proving
\eqref{eq:torsion-S6}.

Fix such a nontrivial $\zeta$, put $K=\Q(\zeta)$, and let $L$ be the
splitting field over $K$.  Since the Galois group of
$L\Q^{\mathrm{ab}}/\Q^{\mathrm{ab}}$ is $S_6$, the group
\[
 \Gal(L/L\cap\Q^{\mathrm{ab}})
\]
has order $720$.  It is a subgroup of $\Gal(L/K)\le S_6$; hence
$\Gal(L/K)=S_6$ and $L\cap\Q^{\mathrm{ab}}=K$.

Put $D=L\cap K^{\mathrm{ab}}$.  Then $D/K$ is a Galois abelian
subextension of the $S_6$-extension $L/K$.  Its corresponding kernel
therefore contains $[S_6,S_6]=A_6$, so
\[
 D\subseteq L^{A_6}.
\]
The discriminant formula \eqref{eq:generic-discriminant} identifies
$L^{A_6}$ with $K(\sqrt{Q(\zeta)})$.  This quadratic extension is
abelian over $K$ and hence is contained in $D$.  Thus equality holds,
which proves \eqref{eq:sign-intersection}; the complementary Galois
group over $K^{\mathrm{ab}}$ is $A_6$, proving
\eqref{eq:A6-over-Kab}.
\end{proof}
